% !TeX root = heights-and-arakelov-geometry.tex
\section{Heights and Arakelov geometry}

The main references are \cite{Yuan-Guo-2025-FirstCourseArakelov}.

We also refer to \cite[Appendices~A.5--A.6 and \S 5.4]{Yuan-Zhang-2026-AdelicLineBundles} and \cite[Chapter~3]{Chambert-Loir-2012-ArakelovGeometryHeights}.

We will use the language of Berkovich space, for reference, see \cite{Berkovich-2012-SpectralTheoryAnalytic,Ducros-Favre-Nicaise-2015-BerkovichSpacesApplications}.

\begin{notation}[label=\labelNFM{not:berkovich-space-and-a-geometry}, status=draft, tag=0000000]
	We fix some notations that will be frequently used in this section.
	\begin{itemize}
		\item A \emph{valuation field} $(\kk,|\cdot|)$ is a field with a (multiplicative) absolute value:
		      $|a| = 0$ if and only if $a = 0$, $|ab| = |a||b|$, and $|a+b| \leq |a|+|b|$.
		      Completeness always refers to the metric $|a-b|$.
		      It is non-archimedean if $|a + b| \leq \max\{|a|,|b|\}$.
		      Two nontrivial absolute values define the same \emph{place} if they define the same topology on $\kk$.
		      For a non-archimedean valuation field, put
		      \[ \calO_\kk = \kk^\circ \coloneqq \{a \in \kk : |a| \leq 1\}, \quad \frakm_\kk = \kk^{\circ\circ} \coloneqq \{a \in \kk : |a| < 1\}, \quad \calk_\kk = \widetilde{\kk} \coloneqq \kk^\circ/\kk^{\circ\circ}. \]
		      Then $\kk^\circ$ is a valuation ring with fraction field $\kk$, unique maximal ideal $\frakm_\kk$, and residue field $\calk_\kk$.
		\item Let $K$ be a number field, $M_K$ its set of places, and $M_K^f$, $M_K^\infty$ its subsets of finite and infinite places.
		      Write $K_v$ for the completion at $v$.
		      For $v \mid p$, normalize $|\cdot|_v$ to extend the usual absolute value on $\bbQ_p$, with $|p|_p = p^{-1}$ at finite places and $\bbQ_\infty = \bbR$.
		      Set $\latexdeletedmath{n_v}\latexaddedmath{\epsilon_v} \coloneqq [K_v:\bbQ_p]/[K:\bbQ]$ for $v \mid p$ and $p \in M_\bbQ$.
		      Thus $\latexdeletedmath{n_v}\latexaddedmath{\epsilon_v} = 1/[K:\bbQ]$ at a real place and $\latexdeletedmath{n_v}\latexaddedmath{\epsilon_v} = 2/[K:\bbQ]$ at a complex place, using the usual complex modulus.
		      Then we have the product formula
		      \[ \sum_{v \in M_K}\latexdeletedmath{n_v}\latexaddedmath{\epsilon_v}\log|a|_v = 0 \quad (a \in K^*). \]
		      Write $\bbC_v$ for the completion of an algebraic closure of $K_v$; at an infinite place this means $\bbC$.
		\item Unless otherwise stated, $X$ is a geometrically integral projective variety over $K$, $d = \dim X \geq 1$, and $L$ is a line bundle on $X$.
		      If $\kk = K$ is a number field, write $X_{v}^\an \coloneqq (X_{K_v})^\an$ for its Berkovich analytification and $X_{\bar{v}}^\an \coloneqq (X_{\bbC_v})^\an$.
		      Also write $X(\bbC_v)$ be the points $\Spec \bbC_v \to X \to \Spec K$ such that the induced $K \to \bbC_v$ are extension of valuation fields ($K$ has valuation $v$).
		      In particular, $X(\bbC_v)$ is disjoint of $[K_v:\bbQ_p]$ copies of $X_{\bbC_v}(\bbC_v)$, here the base change is for a fixed embedding $K \to \bbC_v$.
		      \Yang{}.

		\item Let $\kk$ be a field and $R$ be a ring such that $\Frac R = \kk$.
		      Let $X$ be a projective variety over $\kk$.
		      An \emph{integral model} of $X$ over $R$ is a projective, flat, separated $R$-scheme $\calX$ of finite type, together with a specified isomorphism $\calX_\kk \cong X$.
		      On this section, we only consider the projective case.
		      % If $X/\kk$ is proper (resp. projective), we require that $\calX/R$ is also proper (resp. projective).
		      \Yang{}
	\end{itemize}
	\Yang{}
\end{notation}

\begin{proposition}[label=\labelNFM{prop:spreading-out-for-berkovich}, status=draft, tag=0000000]
	Let $R$ be a domain with fraction field $\kk$, and let $X/\kk$ be a projective variety.
	\begin{enumerate}
		\item There exists $U \subset \Spec R$ non-empty such that $X$ has an integral model over $U$.
		      % If $X$ is proper (resp. projective), the model can be chosen proper (resp. projective).
		\item If $R = \kk^\circ$ with $\kk$ a non-archimedean valuation field, we can take $U = \Spec R$.
		\item Let $\calX,\calX'$ be two integral model of $X/\kk$ over $R$, then there exist a common dominant model $\calX''$ of $X/\kk$ over $R$ with $R$-morphisms $\calX'' \to \calX$, $\calX'' \to \calX'$.
	\end{enumerate}
	\Yang{}
\end{proposition}
\begin{proof}

	\Yang{}
	For spreading out and generic flatness, see \cite[Tags \href{https://stacks.math.columbia.edu/tag/01ZM}{01ZM}, \href{https://stacks.math.columbia.edu/tag/051S}{051S}]{Stacksprojectauthors-2026-StacksProject}.
	For the valuation-ring argument and proper extension of points, see Tags \href{https://stacks.math.columbia.edu/tag/0539}{0539} and \href{https://stacks.math.columbia.edu/tag/0BX4}{0BX4} of the same reference.
	\Yang{}
	% \begin{latexmathchange}
	% Finite presentation allows the equations and gluing data of $X$ to descend to $R[1/s]$.
	% After enlarging $s$, properness or a chosen projective embedding also descends.
	% Generic freeness on a finite affine covering then makes this model flat after a further localization.
	% This explains why the assertion over a general domain is a statement over an open part of its spectrum.

	% Now let $R = \kk^\circ$.
	% If the valuation is trivial, take $\calX = X$.
	% Otherwise choose $a$ with $0 < |a| < 1$.
	% For every $b \in \kk$, one has $a^n b \in R$ for sufficiently large $n$, so $\kk = R[1/a]$.
	% Consequently $X$, viewed as an $R$-scheme, is already a flat model of finite presentation, though its special fiber is empty.
	% If $X$ is proper over $\kk$, apply Nagata compactification over $R$ and take the schematic closure of $X$ in the compactification.
	% Its generic fiber is $X$, because a proper open subscheme of the separated generic fiber is also closed there.
	% The closure is integral, hence torsion free over $R$, hence flat by Tag \href{https://stacks.math.columbia.edu/tag/0539}{0539}.
	% For a projective $X$, use instead its schematic closure in $\bbP^N_R$.
	% Nagata compactification is \cite[Tag \href{https://stacks.math.columbia.edu/tag/0F41}{0F41}]{Stacksprojectauthors-2026-StacksProject}.

	% Finally, for $x \in X^\an$, properness extends $\Spec\scrH(x) \to X$ uniquely to $\Spec\scrH(x)^\circ \to \calX$.
	% The image of the closed point lies in $\calX_\calk$ and defines $\frakr_\calX(x)$.
	% For $x \in X(\overline{\kk})$, completeness gives the unique extension of the absolute value to its finite field of definition.
	% The same valuative criterion gives its reduction.
	% The target is the underlying special-fiber space; its points need not be $\calk$-rational.
	% \end{latexmathchange}
\end{proof}















\subsection[status=draft,tag=0000000,label=\labelNFM{subsec:adic-metric-and-model-metric}]{Berkovich space, adic metric and model metric}

On this subsection, fix a non-archimedean completed valuation field $\kk$.

Let $X$ be a variety over $\kk$.
Recall its Berkovich analytification.

\begin{construction}[label=\labelNFM{constr:berkovich-analytification}, status=draft, tag=0000000]
	Let $X$ be a variety over a complete valuation field $(\kk,|\cdot|)$.
	For an affine open $U = \Spec \calO_X(U) \subset X$, its \emph{Berkovich analytification} is defined as
	\[ U^\an \coloneqq \left\{  \text{  multiplicative seminorm } \rho: \calO_X(U) \to \bbR_{\geq 0} \text{ whose restriction on } \kk \text{ is } |\cdot| \right\}. \]
	We have a natural \emph{forget map}:
	\[ \frakf: U^\an \to |U|, \quad \rho \mapsto \rho^{-1}(0) \in \Spec \calO_X(U). \]
	For every $x \in |U|$ corresponding to prime ideal $\wp_x \subset \calO_X(U)$, we have that $\frakf^{-1}(x)$ is the set of multiplicative norm on $\calO_X(U)/\wp_x$ extending the one on $\kk$.
	And this is naturally bijective the absolute value on $\rkk(x) = \Frac(\calO_X(U)/\wp_x)$ extending the one on $\kk$.
	Hence, we get
	\[ U^\an = \left\{ (x,\gamma) \;\middle|\; x \in |U|, \gamma \text{ an absolute value on } \rkk(x) \text{ extending } |\cdot| \text{ on } \kk \right\}. \]
	Glue these affine pieces, we get a set
	\[ X^\an = \left\{ (x,\gamma) \;\middle|\; x \in |X|, \gamma \text{ an absolute value on } \rkk(x) \text{ extending } |\cdot| \text{ on } \kk \right\}. \]
	Its topology is given by the weak topology with respect to the following maps
	\begin{itemize}[nosep]
		\item the forget map $\frakf: X^\an \to |X|, (x,\gamma) \mapsto x$, where $|X|$ is equipped with the Zariski topology;
		\item For any affine $U \subset X^{\an}$, and any $f \in \calO_X(U)$, the map given by evaluation:
		      \[ ev_f: U^\an \to \bbR_{\geq 0}, (x,|\cdot|_\gamma) \mapsto |\bar{f}|_\gamma, \quad \bar{f} \in \calO_X(U)/\wp_x \subset \rkk(x). \]
		      \Yang{}
	\end{itemize}
	For our now propose, we only consider $X^\an$ as a topological space.
	For more structure, see \cite{Berkovich-2012-SpectralTheoryAnalytic,Ducros-Favre-Nicaise-2015-BerkovichSpacesApplications}.
\end{construction}

There are some useful construction and properties related to the Berkovich analytification.

\begin{construction}[label=\labelNFM{constr:related-construction-in-berkovich-space}, status=draft, tag=0000000]
	Let us recall some construction; see \cite[Section 9.1.1]{Yuan-Guo-2025-FirstCourseArakelov}.
	\begin{itemize}
		\item \emph{Completed residue field.} Let $X/\kk$ be a variety and $(x,\rho) \in X^\an$.
		      The \emph{completed residue field} at $x$ is defined as the completion of $\rkk(x)$ with respect to $\rho$, denoted by $\scrH(x)$.
		      It is an extension of $\kk$ as complete valuation fields.
		\item \emph{Classical points.} Recall that for a finite extension $\bfl/\kk$ with $\kk$ completed valuation field, there exists a unique absolute valuation $\rho_{\bfl/\kk}$ on $\bfl$ extending the one on $\kk$.
		      Then we have a map
		      \[ \fraki: X(\overline{\kk}) \to X^\an, \qquad x \mapsto (x,\rho_{\rkk(x)/\kk}). \]
		      It factor as $X(\overline{\kk}) \to X(\overline{\kk})/G_\kk = \{\text{cloesd points in } |X|\} \injmap X^\an$.
		\item \emph{Reduction map.} Let $X/\kk$ be a projective variety with a projective integral model $\calX/\calO_\kk$.
		      For every $x = (x,\rho) \in X^\an$, by the valuation criterion, we have the diagram
		      \[
			      \begin{tikzcd}
				      \Spec \scrH(x) & \Spec \rkk(x) & \calX \\
				      \Spec \calO_{\scrH(x)} = \{\scrH(x), \calk_{\scrH(x)}\} & & \Spec \calO_\kk = \{\kk,\calk\}.
				      \arrow[from=1-1, to=1-2]
				      \arrow[from=1-2, to=1-3, "x"]
				      \arrow[from=1-1, to=2-1]
				      \arrow[from=1-3, to=2-3]
				      \arrow[from=2-1, to=2-3]
				      \arrow[from=2-1, to=1-3, gray, "\exists! \tilde{x}"]
			      \end{tikzcd}
		      \]
		      Set the reduction map as
		      \[ \frakr_\calX: X^\an \to \calX_\calk, \qquad x \mapsto \tilde{x}(\calk_{\scrH(x)}). \]
		      This reduction map $\frakr_\calX$ extends the reduction map of classical points in \cref{prop:spreading-out-for-berkovich} along $\fraki$.
	\end{itemize}
\end{construction}

\begin{proposition}[label=\labelNFM{prop:properties-of-berkovich-analytification}, status=draft, tag=0000000]
	Let $X,Y$ be varieties over $\kk$ with $\kk$ complete and non-archimedean.
	\begin{enumerate}
		\item A morphism $g:X \to Y$ induces a continuous (indeed analytic) $g^\an:X^\an \to Y^\an$.
		\item The analytification $X^\an$ is compact (resp. connected) iff $X$ is proper over $\kk$ (resp. connected).
		\item Let $\kkk$ be the completion of $\overline{\kk}$, we have $X^\an \cong X_\kkk^\an/G_\kk$ as topological spaces.
	\end{enumerate}
	\Yang{ref}
	% \begin{enumerate}
	% \item \emph{Functoriality} \latexdeleted{and commutative with base change}\latexgrammar{and compatibility with base change.}
	% \begin{latexmathchange}
	% A morphism $g:X \to Y$ induces a continuous $g^\an:X^\an \to Y^\an$, with
	% \[ (h\circ g)^\an = h^\an\circ g^\an, \qquad (\id_X)^\an = \id_{X^\an}. \]
	% For every complete isometric extension $\kk'/\kk$,
	% \[ (X_{\kk'})^\an \cong X^\an\widehat\otimes_\kk\kk'. \]
	% Here $\widehat\otimes$ denotes analytic base change.
	% \end{latexmathchange}
	% \item \latexdeleted{gaga, in particuler, compact}\latexgrammar{Properness and GAGA.}
	% \begin{latexmathchange}
	% If $X$ is projective, $X^\an$ is compact and Hausdorff.
	% With its analytic structure, analytification gives
	% \[ \operatorname{Coh}(X) \simeq \operatorname{Coh}(X^\an), \qquad H^i(X,\scrF) \cong H^i(X^\an,\scrF^\an) \quad(i \geq 0). \]
	% In particular, $H^0(X,L) \cong H^0(X^\an,L^\an)$ for every line bundle $L$.
	% \end{latexmathchange}
	% \end{enumerate}
	% \Yang{}
	% \latexdeleted{To Ai: fill}
\end{proposition}
% \Needspace{14\baselineskip}
% \begin{proof}[References for proof]
% \leavevmode\par
% \begin{latexmathchange}
% On affine charts, $g^\an$ is restriction of seminorms along the map of coordinate rings; this gives functoriality.
% Extension of scalars gives the stated analytic base change.
% The closed unit polydiscs
% \[ \{|X_j/X_i| \leq 1\text{ for all }j\} \subset \bbP^{N,\an}_\kk \quad(0 \leq i \leq N) \]
% are compact and cover projective space.
% A projective $X$ is a closed analytic subspace, hence compact; separatedness gives Hausdorffness.
% The coherent-sheaf and cohomology comparisons are Berkovich GAGA; see \cite[\S\S 3.4--3.5]{Berkovich-2012-SpectralTheoryAnalytic}.
% \end{latexmathchange}
% \end{proof}

Let us turn to the metric on line bundles.

\begin{definition}[label=\labelNFM{def:adic-metric}, status=draft, tag=0000000]
	Let $X$ be a variety over $\kk$ and $L$ a line bundle on $X$.
	A \emph{continuous metric} on $L$ is a family of norms $\|\cdot\| = \{\|\cdot\|_{(x,\gamma)}\}$ such that
	\begin{itemize}
		\item $\|\cdot\|_{(x,\gamma)}$ is a norm on the $(\rkk(x),\gamma)$-vector space $L\otimes \rkk(x)$ for every $(x,\gamma) \in X^\an$;
		\item for any local section $s \in \Gamma(L,U)$, the function $U^\an \to \bbR, (x,\gamma) \mapsto \|s(x)\|_{(x,\gamma)}$ is continuous.
	\end{itemize}
	If $\kk = K_v$ with $K$ number field and $v \in M_K$, we also say $\|\cdot\|$ a \emph{continuous $v$-adic metric}.
	The pair $\bar L = (L,\|\cdot\|)$ is called a \emph{metrized line bundle}.
\end{definition}

As the archimedean case, we have the following definition.

\begin{lemma}[label=\labelNFM{lem:distance-of-two-non-archimead-metrics}, status=draft, tag=0000000]
	Let $X$ be a projective variety over $\kk$ and $L$ a line bundle on $X$.
	All continuous metric on $L$ forms a complete metric space with distance
	\[ \dist(\|\cdot\|_1, \|\cdot\|_2) \coloneqq \max_{x \in X^\an} \left| \log \frac{\|s(x)\|_1}{\|s(x)\|_2} \right|. \]
\end{lemma}
\begin{proof}
	As the archimedean case, all continuous metrics on $L$ is a torsor of $\scrC(X^\an,\bbR)$.
	The infinity norm on $\scrC(X^\an,\bbR)$ gives this metrics and it is well-defined and complete since $X^\an$ is compact by GAGA.
\end{proof}


\begin{example}[label=\labelNFM{eg:standard-metrics-on-projective-space}, status=draft, tag=0000000]
	Let $\bbP^d$ with coordinates $[x_0:\cdots:x_d]$ and $L = \calO(1)$.
	The following rule gives a continuous metric on $L$:
	\[ X_k \mapsto \frac{|X_k|_v}{\max_{0 \leq i \leq d}|X_i|_v}. \]
	This is called the \emph{standard metric} on $\calO(1)$.
\end{example}

Now we introduce a class of metric, which miraculously serves as a non-archimedean analoge of $\cinf$-hermitian metrics.

First let see the complex case, $(X,L)/\bbC$.
To determine a $\cinf$-hermitian metric $\|\cdot\|$ on $L$, we only need to choose an everywhere nonvanish $\cinf$-section $s$ with $\|s(x)\| = 1$ for all $x \in X$.
So, choose an everywhere nonvanish $\cinf$-section $s$, we can get a hermitian metric on $L$ by $\|as\| \coloneqq |a|$ for any $\cinf$-function $a$.
More generally, a $\cinf$-orthogonal frame, i.e., a collection of trivialization $\{(U_i,L|_{U_i} \cong \cinf_{X,\bbC} \cdot e_i)\}$ with transition function $e_i = f_{ij} e_j$ such that $|f_{ij}|=1$, determines a $\cinf$-hermitian metric on $L$ by $\|ae_i\| \coloneqq |a|$ for any $\cinf$-function $a$.
Hence, a $\cinf$-hermitian metric is determined by a collection of transition functions $\{f_{ij}\}$, where
\[ \bigl( f_{ij}: U_{ij} \to S^1 = \{|z|=1\} \subset \bbC^* \bigr) \in \cinf_{X,S^1}(U_{ij}). \]
That is, an element in $H^1(X, \cinf_{X,S^1})$ that has image $[L] \in H^1(X,\calO_X^*) \subset H^1(X,\cinf_{X,\bbC^*})$.

On the non-archimedean place, we have no good $\cinf$-functions.
However, the elements of norm $1$ has some very good structure, namely, we can consider $(\kk^\circ)^*$, the multiplicative group of the ring of integers.
Hence, the question becomes, what is a suitable function class (?).
\[ \bigl( f_{ij}: U_{ij} \to (\kk^\circ)^* = \{|a|=1\} \subset \kk^*\bigr) \in (?). \]
For a scheme $\calY/(\kk^\circ)$, there is a natural class of functions $\calY \to (\kk^\circ)^*$ given by $\calO_\calY^*$.
Therefore, for the variety $X/\kk$, the natural way should be that taking an integral model $\calX/(\kk^\circ)$ and by $X \to \calX \to (\kk^\circ)^*$.
In this way, such a metric should be given by an element in $H^1(\calX,\calO_\calX^*)$ with image $[L] \in H^1(X,\calO_X^*)$.

\begin{construction}[label=\labelNFM{constr:model-metric}, status=draft, tag=0000000]
	Let $X$ be a projective variety over $\kk$ and $L$ line bundle on $X$.
	Given an integral model $(\calX,\calL)$, where $\calX$ is an integral model of $X$, $\calL$ a line bundle on $\latexdeletedmath{X}\latexaddedmath{\calX}$ with an isomorphism $\calL|_X \cong L$,
	then we can define a \emph{model metric} $\|\cdot\|_{(\calX,\calL)}$ by the frame of $\calL$ on $X$.
	More explicitly, let $\{(\calU_i,\calL|_{\calU_i} = \calO_{\calU_i} \cdot e_i) \}$ be a trivialization of $\calL$.
	\latexdeleted{This also gives a trivialization}
	\[ \latexdeletedmath{\{((\calU_i \cap X \eqqcolon U_i)^{\an}, L|_{U_i}^{\an} = \calL|_{U_i}^{\an} \cong \scrO_{U_i^\an} \cdot e_i\}} \]
	\latexdeleted{of $L^\an$ on $X^\an$.}
	\latexdeleted{For any local continuous section $ae_i$, set $\|ae_i(x)\|_{(\calX,\calL)}\coloneqq|a|_x$.}
	\begin{latexmathchange}
		The unit-frame rule is used on reduction tubes:
		\[ V_i \coloneqq \frakr_\calX^{-1}(\calU_i\cap\calX_\calk), \qquad \|ae_i(x)\|_{(\calX,\calL)} \coloneqq |a(x)|\quad(x \in V_i). \]
		The $V_i$ are analytic domains covering $X^\an$.
		On $V_i\cap V_j$, the transition $u_{ij}$ and its inverse are integral at the reduction of $x$, hence
		\[ |u_{ij}(x)| \leq 1, \qquad |u_{ij}^{-1}(x)| \leq 1, \qquad |u_{ij}(x)| = 1. \]
		Thus the local rules agree and give a continuous metric.
	\end{latexmathchange}

	For latter use, we also need a version with rational coefficients.
	Given an integral model $(\calX,\calM)$ of $(X,L^{\otimes m})$, the \emph{model metric (with rational coefficients)} of $(\calX,\calM, m)$ is defined as
	\[ \|s\|_{(\calX,\calM,m)} \coloneqq \|s^{\otimes m}\|_{(\calX,\calM)}^{1/m}.\]
\end{construction}

\Needspace{16\baselineskip}
\begin{remark}[label=\labelNFM{rmk:generic-fiber-and-reduction-domains}, status=draft, tag=0000000]
	\leavevmode\par
	\begin{latexmathchange}
		The frame $e_i$ does trivialize $L^\an$ on $U_i^\an$, where $U_i = \calU_i\cap X$.
		The smaller domain $V_i \subset U_i^\an$ requires the reduction of $x$ to lie in $\calU_i\cap\calX_\calk$.
		For example, take $(\calX,\calL) = (\bbP^1_{\kk^\circ},\calO(1))$, with
		\[ \calU_0 = \{X_0 \ne 0\}, \qquad \calU_1 = \{X_1 \ne 0\}, \qquad t = X_1/X_0. \]
		Then
		\[ U_0^\an = (\bbA^1_\kk)^\an, \qquad V_0 = \{x \in U_0^\an : |t(x)| \leq 1\}. \]
		Choose $\pi \in \kk$ with $0 < |\pi| < 1$.
		The point $x = [1:\pi^{-1}]$ belongs to $U_0^\an$, but
		\[ x = [\pi:1] \quad\Longrightarrow\quad \frakr_\calX(x) = [0:1] \notin \calU_0\cap\calX_\calk, \qquad x \notin V_0. \]
		Moreover,
		\[ U_0^\an\cap U_1^\an = (\bbG_{m,\kk})^\an, \qquad V_0\cap V_1 = \{x \in (\bbG_{m,\kk})^\an : |t(x)| = 1\}. \]
		The frames $e_0 = X_0$, $e_1 = X_1$ satisfy $e_1 = te_0$.
		Although $t \in \calO_\calX(\calU_0\cap\calU_1)^*$, we have $|t(x)| = |\pi|^{-1} > 1$ at the above point.
		Thus the prescriptions $\|e_0\| = \|e_1\| = 1$ agree on $V_0\cap V_1$, but not on the whole $U_0^\an\cap U_1^\an$.
	\end{latexmathchange}
\end{remark}

\begin{example}[label=\labelNFM{eg:model-metric-by-standard-p1-and-blowup}, status=draft, tag=0000000]
	Let $X = \bbP^1_\kk$ and $L = \calO_X$.



	Consider the standard model $(\bbP^d_{\calO_\kk},\calO(1))$ of $(\bbP^d_{\kk},\calO(1))$.


	Let $\|\cdot\|$ be the standard metric on $(\bbP^d,\calO(1))$.

	\Yang{}
\end{example}

\begin{example}[label=\labelNFM{eg:standard-metric-are-model-metric}, status=draft, tag=0000000]
	Consider the standard model $(\bbP^d_{\calO_\kk},\calO(1))$ of $(\bbP^d_{\kk},\calO(1))$.


	Let $\|\cdot\|$ be the standard metric on $(\bbP^d,\calO(1))$.

	\Yang{}
\end{example}

Recall that in the archimedean case, any $\cinf_{-,\bbR}$-function can be got by $-\log \|1\|$ for some $\cinf$-metric $\|\cdot\|$ on $\scrO_X$.
By this way, we have the following define.

\begin{definition}[label=\labelNFM{def:model-function-over-na-place}, status=draft, tag=0000000]
	Let $X$ be a projective variety over $\kk$.
	A \emph{model function} on $X$ is a function $X^\an \to \bbR$ of form $- \log \|1\|_{(\calX,\calL,m)}$, where $(\calX,\calL,m)$ is a model of $(X,\calO_X)$ and $1 \in \calO_X(X)$ is the canonical section.
\end{definition}

\Needspace{18\baselineskip}
\begin{example}[label=\labelNFM{eg:model-function}, status=draft, tag=0000000]
	\latexdeleted{To Ai: give some example and nonexample of model function}
	\begin{latexmathchange}
		Let $\pi$ be a uniformizer of $K_v$.
		Since the special fiber of a model is the divisor of $\pi$, every rational multiple of $-\log|\pi|_v$ is a constant model function.

		For a nonconstant example, let $X = \bbP^1_K$ with affine coordinate $z$ and put
		\[ \varphi(z) \coloneqq \log\frac{\max\{|\pi|_v,|z|\}}{\max\{1,|z|\}}. \]
		The expression extends continuously to $\bbP^{1,\an}_{\bbC_v}$, with value $0$ at $\infty$.
		It is a model function: on $\bbP^1_{K_v^\circ}$, take the vertical ideal $(\pi,z)$ on the chart $z \neq \infty$ and the unit ideal on the other chart.
		Its blowup makes the ideal invertible; the negative of the exceptional divisor gives $\varphi$.

		The continuous function $\varphi^2$ is \emph{not} a model function.
		Indeed, on the segment of Gauss points $x_{0,r}$ with $|\pi|_v \leq r \leq 1$,
		\[ \varphi(x_{0,r}) = \log r, \qquad \varphi^2(x_{0,r}) = (\log r)^2. \]
		A rational model function restricts to a piecewise affine function of $\log r$ with finitely many pieces,
		whereas $(\log r)^2$ does not.
		For the piecewise affine behavior on skeleta, see \href{https://gubler.app.uni-regensburg.de/Publikationen/1404.7044v1.pdf}{Gubler--Rabinoff--Werner}.
	\end{latexmathchange}
\end{example}

\begin{proposition}[label=\labelNFM{prop:basic-properties-of-model-function}, status=draft, tag=0000000]
	Let $X/\kk$ be a projective variety.
	, with $\kk$ complete, nontrivially valued and non-archimedean.
	Write $\scrC^\rmM(X)$ for the rational model functions.
	\begin{enumerate}
		\item For $f,g \in \scrC^\rmM(X)$ and $a,b \in \bbQ$, $af+bg \in \scrC^\rmM(X)$.
		\item If $f,g \in \scrC^\rmM(X)$, then $\max\{f,g\},\min\{f,g\} \in \scrC^\rmM(X)$.
		\item The $\bbQ$-vector space $\scrC^\rmM(X)$ is dense in $(\scrC(X^\an,\bbR), \|\cdot\|_\infty)$.
	\end{enumerate}
	\Yang{}
\end{proposition}
% \begin{proof}[References for proof]
% \leavevmode\par
% \Yang{}
% \begin{latexmathchange}
% On a common model, write $f = \varphi_D$, $g = \varphi_E$ for vertical rational Cartier divisors, with $\varphi_D = -\log\|1\|_{\calO(D)}$.
% Then $af+bg = \varphi_{aD+bE}$.
% After clearing denominators, blow up the sum of the corresponding vertical fractional ideals.
% Its invertible pullback induces $\max\{f,g\}$; use $\min\{f,g\} = -\max\{-f,-g\}$.
% See \cite[Proposition~3.12]{Gubler-Martin-2019-ZhangSemipositiveMetrics}.
% \end{latexmathchange}
% \begin{latexmathchange}
% For $L = \calO_X$, uniform approximation of a continuous metric by piecewise $\bbQ$-linear metrics gives the density statement.
% On proper $X$, these are precisely rational model metrics; see \cite[Theorems~1.1 and~2.17]{Gubler-Martin-2019-ZhangSemipositiveMetrics}.
% Nontriviality of the valuation is needed here: it supplies the dense group of constant model functions $\bbQ\log|a|$, for $0 < |a| < 1$.
% \end{latexmathchange}
% \end{proof}

\Needspace{16\baselineskip}
\begin{remark}[label=\labelNFM{rmk:tropical-semiring-and-model-function}, status=draft, tag=0000000]
	\latexdeleted{The model functions usual does not form a ring since it cannot take multiplication.}
	\latexdeleted{but To Ai: fill this remark}
	\begin{latexmathchange}
		Model functions form a $\bbQ$-vector space and are closed under pointwise maximum and minimum.
		For sums and rational multiples, pass to a common model and add or scale the vertical divisors.
		To form a maximum, after passing to a common model, blow up the sum of the corresponding vertical fractional ideals so that it becomes invertible.
		The resulting model function is the pointwise maximum; the minimum follows by taking negatives.

		Define tropical operations by
		\[ \varphi \oplus \psi \coloneqq \max\{\varphi,\psi\}, \qquad \varphi \odot \psi \coloneqq \varphi + \psi. \]
		After adjoining the constant function $-\infty$, these operations make $\calM_v(X) \cup \{-\infty\}$ an idempotent semiring,
		with tropical zero $-\infty$ and tropical unit the ordinary zero function.
		Ordinary pointwise products need not be model functions, as $\varphi^2$ in the preceding example shows.
	\end{latexmathchange}
\end{remark}


By the model metric, we can define the positivity of non-archimedean continuous metrics.

\begin{definition}[label=\labelNFM{def:semipositive-local-metric}, status=draft, tag=0000000]
	Let $X$ be a projective variety over $\kk$ and $L$ a line bundle on $X$.
	A model metric $\|\cdot\|_{(\calX,\calL,m)}$ is called \emph{semipositive} if $\calL$ is relatively nef, that is, $\deg(\calL|_C) \geq 0$ for every integral curve $C$ in the special fiber.
	A continuous metric is called \emph{semipositive} if it is a uniform limit of semipositive model metrics.
\end{definition}

\begin{lemma}[label=\labelNFM{lem:semipositive-is-well-defined}, status=draft, tag=0000000]
	Let $X$ be a projective variety over $\kk$, $L$ a line bundle on $X$ and $\|\cdot\|$ a model metric on $L$.
	Then $\|\cdot\|$ is semipositive as model metric iff as continuous metric.
	\Yang{}
\end{lemma}

\begin{definition}[label=\labelNFM{def:Chambert-Loir-measure}, status=draft, tag=0000000]
	Let $X$ and
	Let $\bar L$ be nef, with $L$ big.
	At every place $v$, there is a positive measure $c_1(\bar L_v)^d$ on $X_v^\an$ of total mass $\deg_L(X)$.
	At an infinite place it is the Bedford--Taylor product of the curvature current.
	At a finite place it is the Chambert-Loir measure: for model metrics it is obtained from intersection degrees on the special fiber, and semipositive metrics are treated by weak limits.
	We set
	\[ \mu_{\bar L,v} \coloneqq \frac{c_1(\bar L_v)^d}{\deg_L(X)}. \]
	Thus $\mu_{\bar L,v}$ is a probability measure.
	\Yang{}
\end{definition}




























\subsection[status=draft,tag=0000000,label=\labelNFM{subsec:adelic-line-bundle-and-positivity}]{Adelic line bundle and arithmetic intersection}

On this subsection, fix a number field $K$.

\begin{definition}[label=\labelNFM{def:adelic-line-bundle-projective-case}, status=draft, tag=0000000]
	Let $X$ be a projective variety over $K$.
	An \emph{adelic metric} on $X$ is a collection

	An \emph{adelic line bundle} on $X$ is a pair
	\[ \bar L = (L,(\|\cdot\|_v)_{v \in M_K}) \]
	with continuous $v$-adic metrics, such that there is an nonempty open subset $U \subset \Spec \calO_K$ and an integral model $(\calX,\calM,m)/U$ of $(X,L)$ such that for all $v \in U$, $\|\cdot\|_v = \|\cdot\|_{(\calX|_{\calO_v},\calM|_{\calO_v},m)}$ is the model metric induced by the pullback of $(\calX,\calL)$.
	Their isometry classes form $\aPic(X)$ under the local operators.

	If we can take $U = \Spec \calO_K$ and $\|\cdot\|_v$ being $\cinf$ for $v \in M_K^\infty$, then we say that $\overline{L} \eqqcolon \overline{\calL} = (\calX,\calM,m; (\|\cdot\|_v)_{v \in M_K^\infty})$ is a \emph{$\cinf$--model metric line bundle}.
\end{definition}
A morphism $g:Y \to X$ gives a pullback $g^*:\aPic(X) \to \aPic(Y)$.

\begin{construction}[label=\labelNFM{constr:topology-on-the-adelic-metrics-on-a-line-bundle}, status=draft, tag=0000000]
	Let $X/K$ be a projective variety and $L$ a line bundle on $X$.
	In the adelic case, all adelic metric on $L$ also consists of a topological space.
	First note that any two adelic metric on $L$ only differen on finitely many place: pullback them to the same model $\calX/U$, since $\calL_1|_X \cong \calL_2|_X$, after thrinking $U$, we have $\calL_1|_U \cong \calL_1|_U$.

	Fix a metric $\|\cdot\|$ on $L$.
	For any nonempty open subset $U \subset \Spec \calO_K$,
	\[ E_U \coloneqq \left\{ \|\cdot\|' \text{ adelic metric on } L  \;\middle|\; \|\cdot\|'_v = \|\cdot\|_v, \forall v \in U \right\}. \]
	The space $E_U$ is a complete metric space (in fact a banach space with $\|\cdot\|$ as origin) with distance
	\[  \dist_U(\|\cdot\|',\|\cdot\|'') \coloneqq \sum_{v \in M_K \setminus U}\max_{x \in X^\an_v} \left| \log \frac{\|s(x)\|_v'}{\|s(x)\|_v''} \right|. \]
	For every $V \subset U$, we have continuous inclusion $E_U \subset E_V$.
	Hence we have the inductive limit of sets:
	\[ \{\text{ adelic metrics on } L\} = \dirlim_{U \subset \Spec \calO_K} E_U \]
	On the set of adelic metrics on $L$, we equip it with the inductive limit topology.
	That is, $\|\cdot\|_n \to \|\cdot\|_0$ iff there exists $U$ such that $\|\cdot\|_n \in E_U$ for all $n$ and $\|\cdot\|_n \to \|\cdot\|_0$ in $E_U$.
	Easily see that the topology is independent of choice of $\|\cdot\|$.
\end{construction}

\begin{definition}[label=\labelNFM{def:height-of-an-algebraic-point}, status=draft, tag=0000000]
	Let $x \in X(\bar K)$ and choose a finite extension $E/K$ with $x \in X(E)$.
	For a rational section $s$ of $L$ regular and nonzero at $x$, define
	\[ h_{\bar L}(x) \coloneqq -\sum_{w \in M_E} \latexdeletedmath{n_w}\latexaddedmath{\epsilon_w}\log\|s(x)\|_w. \]
	\latexdeleted{To Ai: change all $n_v$ to $\epsilon_v$.}
\end{definition}

The metrics are pulled back from $K$.
On the sum, only finitely many summands are nonzero.
Replacing $s(x)$ by $a s(x)$ changes the sum by $-\sum_w \latexdeletedmath{n_w}\latexaddedmath{\epsilon_w}\log|a|_w = 0$.
The local degree formula shows that the value is also independent of $E$.
Thus this gives a well-defined real-valued height on $X(\bar K)$.

\begin{example}[label=\labelNFM{eg:standard-logarithmic-height}, status=draft, tag=0000000]
	Give $\scrO(1)$ on $\bbP^d$ the standard maximal metric at every place, and denote the result by $\overline{\scrO(1)}_{\max}$.
	Applying the definition with a nonzero coordinate section (for example, $X_0$,) gives
	\begin{align*}
		h_{\overline{\scrO(1)}_{\max}}([a_0:\cdots:a_d]) & = -\sum_{w \in M_E}\latexdeletedmath{n_w}\latexaddedmath{\epsilon_w}\log |a_0|_w + \sum_{w \in M_E}\latexdeletedmath{n_w}\latexaddedmath{\epsilon_w}\log\max_i|a_i|_w \\
		                                                 & = \sum_{w \in M_E}\latexdeletedmath{n_w}\latexaddedmath{\epsilon_w}\log\max_i|a_i|_w.
	\end{align*}
	The product formula removes the term involving that coordinate.
	In particular, for $x = [x_0:\cdots:x_d] \in \bbP^d(\bbQ)$ with $\gcd(\{x_i\}) = 1$, we have
	\[  h_{\overline{\scrO(1)}_{\max}}(x) = \log \max\{|x_0|,\cdots,|x_d|\}. \]
\end{example}

\begin{proposition}[label=\labelNFM{prop:fundamental-properties-of-heights}, status=draft, tag=0000000]
	We have the following properties:
	\begin{enumerate}
		\item (Linearity) For adelic $\bbQ$-line bundles and $a,b \in \bbQ$,
		      \[ h_{a\bar L+b\bar M}(x) = a h_{\bar L}(x)+b h_{\bar M}(x). \]
		\item (Functoriality) For $g:Y \to X$ and $y \in Y(\bar K)$,
		      \[ h_{g^*\bar L}(y) = h_{\bar L}(g(y)). \]
		\item (Boundedness) Let $\overline{L}_1$,$\overline{L}_2$ be two adelic line bundle with the same underlying line bundle, then there exists $M > 0$ such that
		      \[ | h_{\overline{L}_1}(x) - h_{\overline{L}_2}(x) | \leq M,\qquad \forall x \in X(\overline{K}). \]
		\item (Northcott properties) If $L$ is ample, then for every $D \geq 1$ and $B \in \bbR$,
		      \[ \#\{x \in X(\bar K):[K(x):K] \leq D,\ h_{\bar L}(x) \leq B\} < \infty. \]
	\end{enumerate}
	\Yang{}
\end{proposition}
\Needspace{18\baselineskip}
\begin{proof}
	\leavevmode\par
	\begin{latexmathchange}
		Linearity and functoriality follow by tensoring sections and pulling back their local norms.
		Base extension follows from
		\[ \epsilon_w = \frac{[E_w:K_v]}{[E:K]}\epsilon_v, \qquad \sum_{w\mid v}\epsilon_w = \epsilon_v. \]
		For ample $L$, choose $m > 0$ and a closed immersion $j:X\hookrightarrow\bbP^N$ with $L^{\otimes m} = j^*\calO(1)$.
		Compactness and agreement of the metrics outside finitely many places give
		\[ m h_{\bar L}(x) = h_{\mathrm{std}}(j(x))+O(1), \]
		with a bound independent of $x$ and its field of definition.
		The standard height is nonnegative and has Northcott finiteness; this proves both assertions.
		See \cite[Appendix~A.1]{Yuan-Zhang-2026-AdelicLineBundles}.
	\end{latexmathchange}
\end{proof}

\Yang{}

View $M_K$ as a ``curve'', \Yang{}

\begin{definition}[label=\labelNFM{def:positivity-of-adelic-line-bundles}, status=draft, tag=0000000]
	An adelic line bundle is \emph{semipositive}, or \emph{relatively nef}, if all its local metrics are semipositive.
	It is \emph{nef} if it is semipositive and $h_{\bar L}(x) \geq 0$ for every $x \in X(\bar K)$.
	It is \emph{integrable} if $\bar L \cong \bar A \otimes \bar B^\vee$ for two nef adelic $\bbQ$-line bundles $\bar A,\bar B$.
	We denote the group of integrable adelic line bundles by $\aPic(X)_{\mathrm{int}}$.
	\Yang{to add positive of metric in archimedean place}
\end{definition}

Tensor products, positive rational powers, and pullbacks preserve semipositivity and nefness.

\begin{construction}[label=\labelNFM{constr:intersection-number-of-adelic-line-bundles}, status=draft, tag=0000000]
	Let $X$ be a projective variety over $K$.
	We will try to construction the intersection number of $\cinf$--model adelic line bundles on $X$.

	We begin from hermitian line bundles.
	Let $\overline{\calL}_0,\cdots,\overline{\calL}_d$ be hermitian lines on an integral model $\calX/\calO_K$.
	Choose a nonzero rational section $l_d$ of $\overline{\calL}_d$ and write $\divisor l_d = \sum a_i \calH_i + \sum b_i V_j$ with $\calH_i,V_j$ prime divisors on $\calX$ and $\calH_i$ (resp. $V_j$) is horizontal (resp. vertical) over $\Spec \calO_K$.
	Set $p(V_j)$ be the characteristic of $V_j$, i.e., the image of $V_j$ under $\calX \to \Spec \bbZ$ is $(p)$.
	we inductively define
	\begin{align*}
		  & \overline{\calL}_0 \cdots \overline{\calL}_d                                                                                                                    \\
		= & \sum a_i \cdot \overline{\calL}_0|_{\calH_i} \cdots \overline{\calL}_{d-1}|_{\calH_i} + \sum b_j \cdot \calL_0 \cdots \calL_{d-1} \cdot [V_j] \cdot \log p(V_j) \\
		  & - \sum_{v \in M_K^\infty} \latexaddedmath{[K_v:\bbR]}\int_{\calX(\bbC_v)} \log \|l_d\|_v \cdot c_1(\overline{\calL}_0)\cdots c_1(\overline{\calL}_{d-1}).
	\end{align*}
	The first item is defined inductively.
	The second item is the classical intersection number (over the prime field $\bbF_p$) with weight $\log p$ by its characteristic.
	The last item compute the contribution of archimedean places.
	\begin{latexmathchange}
		The sum is over places; $[K_v:\bbR] = 1$ or $2$ accounts for real embeddings or conjugate pairs.
	\end{latexmathchange}
	Here $\calX(\bbC_v)$ means the points $\Spec \bbC \to \calX \to \Spec \calO_K$ such that $K \to \bbC$ coincides with $v:K \to \bbC$.
	\latexdeleted{When $\dim x = 0$, the first item disappear and the formula coincides with the height.}
	\begin{latexmathchange}
		For $X = \Spec E$, this unnormalized degree is $[E:\bbQ]h_{\bar L}(x)$. Normalized heights divide by the absolute field degree.
	\end{latexmathchange}
	\Yang{}

	The intersection pairing is multilinear and symmetric \Yang{ref}.
	To extend the definition to general adelic line bundle, we may expect some continuity.
	However, this fails in general.
	But it is continuous for semipositive ($\cinf$--model) adelic line bundles \Yang{ref}.
	Hence we can define the intersection pairing for all integrable adelic line bundle by continuation and Linearity.

	More generally, let $Z\subset X$ be a subvariety of dimension $r$, we can define
	\[ \overline{L}_0 \cdots \overline{L}_r \cdot [Z] \coloneqq \overline{L}_0|_Z \cdots \overline{L}_r|_Z. \]
	This extends to $Z_r(X)$ by linearity.
	Final, we can define the height for a subvariety $Z$ with respect to $\overline{L}=(L,\|\cdot\|)$ by
	\[ \latexdeletedmath{h_{\bar L}(Z) \coloneqq \frac{(\bar L)^{r+1}\cdot[Y]}{(r+1)(L)^r\cdot [Y]}.} \]
	\begin{latexmathchange}
		For $\deg_L(Z) = (L^r\cdot[Z]) > 0$, use the normalized height
		\[ h_{\bar L}(Z) \coloneqq \frac{(\bar L)^{r+1}\cdot[Z]}{(r+1)[K:\bbQ]\deg_L(Z)}. \]
		In dimension zero, $(\bar L\cdot[P]) = [K(P):\bbQ]h_{\bar L}(x)$ and $\deg_L(P) = [K(P):K]$, so this agrees with the point height.
	\end{latexmathchange}
	When $z$ is a point, this coincides with usual height.\Yang{}
\end{construction}

Let us list some basic properties of the intersection number of adelic line bundles.

\begin{proposition}[label=\labelNFM{prop:fundamental-properties-of-arithmetic-intersection}, status=draft, tag=0000000]
	\latexdeleted{The}
	\begin{latexmathchange}
		Let $X/K$ be projective, $Z \in Z_r(X)$, and $\bar L_0,\ldots,\bar L_r$ integrable adelic $\bbQ$-line bundles.
		All intersections below are unnormalized.
	\end{latexmathchange}
	\begin{enumerate}
		\item The intersection pairing $\aPic(X)_{int}^{d+1} \to \bbR$ is multilinear and symmetric.
		\item For a finite extension $E/K$, put $X_E = X\times_K E$ and let $Z_E$ be the flat pullback of $Z$.
		      Then
		      \[ (\bar L_{0,E}\cdots\bar L_{r,E}\cdot Z_E)_E = (\bar L_0\cdots\bar L_r\cdot Z)_K. \]
		\item For a morphism $g:Y \to X$ of projective $K$-varieties and $W \in Z_r(Y)$,
		      \[ (g^*\bar L_0\cdots g^*\bar L_r\cdot W) = (\bar L_0\cdots\bar L_r\cdot g_*W). \]
		      See \cite[Proposition~4.1.2 and Lemma~4.6.1(3)]{Yuan-Zhang-2026-AdelicLineBundles}.
		\item
		      Let $\bar L_0'$ be integrable, with the same underlying bundle as $\bar L_0$ and the same metrics outside $v$.
		      If $\|\cdot\|_v' = e^{-\varphi}\|\cdot\|_v$, then
		      \begin{align*}
			       & (\bar L_0'\bar L_1\cdots\bar L_r\cdot Z)-(\bar L_0\bar L_1\cdots\bar L_r\cdot Z)                    \\
			       & \qquad = [K:\bbQ]\epsilon_v\int_{Z_{\bar v}^\an}\varphi\,c_1(\bar L_{1,v})\cdots c_1(\bar L_{r,v}).
		      \end{align*}
		      Here the metrics and $\varphi$ are pulled back to $\bbC_v$, and the integral is linear in $Z$.
		      See \cite[\S 4.2.2, Theorem~4.6.2 and \S A.6.2]{Yuan-Zhang-2026-AdelicLineBundles}.
	\end{enumerate}

	\Yang{}
\end{proposition}

% \begin{theorem}[label=\labelNFM{thm:}, status=draft, tag=0000000]
% 	For integrable adelic $\bbQ$-line bundles, the arithmetic intersection has the following properties.
% 	\begin{enumerate}
% 		\item It is symmetric and multilinear in the line bundles and additive in the cycle.
% 		      It is continuous for uniform convergence of semipositive metrics with fixed underlying bundles and fixed metrics outside a finite set of places.
% 		\item If $g:Y \to X$ is a morphism of projective varieties and $W \subset Y$ is an integral $r$-dimensional subvariety, the projection formula is
% 		      \[ I_K(g^*\bar L_0,\ldots,g^*\bar L_r;W) = I_K(\bar L_0,\ldots,\bar L_r;g_*[W]). \]
% 		      Here $g_*[W] = [K(W):K(g(W))][g(W)]$ when $\dim g(W) = r$, and is zero otherwise.
% 		\item For a finite extension $E/K$, with base-changed cycles interpreted additively,
% 		      \[ I_E(\bar L_{0,E},\ldots,\bar L_{r,E};Z_E) = [E:K]I_K(\bar L_0,\ldots,\bar L_r;Z). \]
% 		      Thus normalized heights are unchanged by base extension.
% 		\item Let $r = \dim X = d$, and change only the metric of $\bar L_0$ at $v$ by $\|\cdot\|_v' = e^{-\varphi}\|\cdot\|_v$, with the changed bundle integrable.
% 		      Then
% 		      \[ I_K(\bar L_0',\bar L_1,\ldots,\bar L_d;X)-I_K(\bar L_0,\bar L_1,\ldots,\bar L_d;X) = [K:\bbQ]\epsilon_v\int_{X_{\bar v}^\an}\varphi\,c_1(\bar L_{1,v})\cdots c_1(\bar L_{d,v}). \]
% 		\item If all the bundles are nef, their intersection on an effective cycle is nonnegative.
% 	\end{enumerate}
% 	See \cite[\S\S 4.1--4.4 and Appendix~A.5.3]{Yuan-Zhang-2026-AdelicLineBundles}.
% 	\Yang{}
% \end{theorem}
% \begin{proof}
% 	\leavevmode\par
% 	\begin{latexmathchange}
% 		For model metrics these are the symmetry, multilinearity, projection formula, metric-change formula, and positivity of arithmetic intersection.
% 		When the other bundles are semipositive, the metric-change formula bounds a variation by the uniform norm of the metric change times the relevant geometric degree.
% 		It therefore gives continuity and permits passage to the defining limits of nef adelic bundles.
% 		Expanding differences of nef bundles gives the integrable case; base change follows either on models or from the local-degree formula for the metrized Deligne pairing.
% 	\end{latexmathchange}
% \end{proof}

The heights of point can be controlled by heights of varieties.
The first version of following is proposed by Shou-Wu Zhang.

\begin{definition}[label=\labelNFM{def:successive-minima}, status=draft, tag=0000000]
	\Yang{}
	Let $\dim X = d$ and $\bar L$ be an adelic line bundle with $L$ ample.
	For $1 \leq i \leq d+1$, define the \emph{successive minima} by
	\[ e_i(\bar L) \coloneqq \sup_{\substack{Z \subset X_{\bar K}\text{ closed}\\\dim Z \leq d-i}}\ \inf_{x \in X(\bar K)\setminus Z(\bar K)} h_{\bar L}(x). \]
	We allow $Z = \varnothing$, with dimension $-\infty$.
\end{definition}

\Yang{}
Then $e_1 \geq \cdots \geq e_{d+1}$; $e_1$ is the \emph{essential minimum}, and $e_{d+1} = \inf_x h_{\bar L}(x)$ is the absolute minimum.
Equivalently,
\[ e_1(\bar L) = \inf\{t \in \bbR:\{x \in X(\bar K):h_{\bar L}(x) \leq t\}\text{ is Zariski dense}\}. \]
For every generic sequence $(x_m)$, one has $\liminf_m h_{\bar L}(x_m) \geq e_1(\bar L)$.
Indeed, it eventually avoids each fixed proper closed $Z$ in the definition.
The same definition of $e_1$ also makes sense when $L$ is only big.
See \cite[Appendix~A.5.5]{Yuan-Zhang-2026-AdelicLineBundles}; that appendix uses arithmetic dimension $d+1$ in place of our geometric dimension $d$.

\begin{theorem}[label=\labelNFM{thm:fundamental-inequality-IV}, status=draft, tag=0000000, title={{{Fundamental inequality \Yang{}}}}]
	\Yang{}
	Let $X/K$ be projective of dimension $d$, let $L$ be ample, and let $\bar L$ be semipositive.
	Zhang's theorem of successive minima states
	\[ e_1(\bar L) \geq h_{\bar L}(X) \geq \frac{e_1(\bar L)+\cdots+e_{d+1}(\bar L)}{d+1}. \]
	If $\bar L$ is nef with big underlying $L$, one still has the fundamental inequalities
	\[ e_1(\bar L) \geq h_{\bar L}(X) \geq \frac{e_1(\bar L)}{d+1}. \]
	For the first statement and proof, see \cite[Theorem~A.5.2]{Yuan-Zhang-2026-AdelicLineBundles}; for the nef and big version, see Theorem~5.3.3 there.
\end{theorem}


As classic case, we can define the volume of adelic line bunle.

\begin{definition}[label=\labelNFM{def:arithmetic-volume-of-adelic-line-bundle}, status=draft, tag=0000000]
	Let $X/K$ be projective of dimension $d$ and $\bar L$ an adelic line bundle.
	The set of \emph{small sections} and its logarithmic cardinality are
	\[ \widehat H^0(X,m\bar L) \coloneqq \{s \in H^0(X,L^{\otimes m}):\|s\|_{v,\sup} \leq 1\text{ for every }v\}, \qquad \widehat h^0(X,m\bar L) \coloneqq \log\#\widehat H^0(X,m\bar L). \]
	The \emph{arithmetic volume} is
	\[ \widehat{\vol}(\bar L) \coloneqq \limsup_{m \to \infty}\frac{\widehat h^0(X,m\bar L)}{m^{d+1}/(d+1)!}. \]
	The bundle is \emph{arithmetically big} if $\widehat{\vol}(\bar L) > 0$.
	\Yang{}
\end{definition}

\Yang{}

This is a finite set containing the zero section, rather than a $K$-vector space.

For $\bbQ$-line bundles, clear denominators and use homogeneity of degree $d+1$.
This is stronger than bigness of the underlying $L$, which concerns growth of $\dim_K H^0(X,mL)$ in degree $d$.
Our volume has no factor $1/[K:\bbQ]$.
For a nef $\bar L$, the arithmetic Hilbert--Samuel formula gives
\[ \widehat{\vol}(\bar L) = \latexdeletedmath{I_K(\bar L,\ldots,\bar L;X)}\latexaddedmath{(\bar L)^{d+1}\cdot[X]}. \]
See \cite[\S\S 5.1--5.2, especially Theorem~5.2.2]{Yuan-Zhang-2026-AdelicLineBundles}.

\Yang{}

\begin{theorem}[label=\labelNFM{thm:arithmetic-hilbert-samuel-siu-inequality}, status=draft, tag=0000000]
	Let $X/K$ be projective of dimension $d$, \latexdeleted{and let}
	\begin{enumerate}
		\item (Arithmetic Hilbert--Samuel formula). If $\bar L$ is nef, then
		      \[ \widehat h^0(X,m\bar L) = \frac{(\bar L)^{d+1}\cdot[X]}{(d+1)!}m^{d+1}+o(m^{d+1}), \]
		      and hence $\widehat{\vol}(\bar L) = (\bar L)^{d+1}\cdot[X]$.
		\item (Arithmetic Siu inequality). If $\bar L,\bar M$ are nef, then
		      \[ \widehat{\vol}(\bar L-\bar M) \geq (\bar L)^{d+1}\cdot[X]-(d+1)(\bar L)^d\bar M\cdot[X]. \]
		      In particular, a positive right-hand side implies that $\bar L-\bar M$ is arithmetically big.
	\end{enumerate}
	\Yang{}
\end{theorem}
\begin{proof}[References for proof]
	\leavevmode\par
	\Yang{}
	\begin{latexmathchange}
		These are \cite[Theorem~5.2.2(1)--(2)]{Yuan-Zhang-2026-AdelicLineBundles}, together with existence of the volume limit in Theorem~5.2.1.
		For hermitian model bundles they are the arithmetic Hilbert--Samuel and arithmetic bigness theorems.
		Nef approximation and continuity of volume and intersection give the adelic statements; see also \S 5.2 there.
		The degree in that reference is the arithmetic dimension, hence our coefficient $d+1$.
	\end{latexmathchange}
\end{proof}


































\subsection[status=draft,tag=0000000,label=\labelNFM{subsec:canonical-height}]{Canonical heights}

On this subsection, let us turn to the dynamics.

\begin{lemma}[label=\labelNFM{lem:existence-of-canonical-adelic-line-bundle}, status=draft, tag=0000000]
	Let $X/K$ be a projective variety and $f: X \to X$ a $q$-polarized endomorphism.
	Fix a polarization $f^*L \cong L^{\otimes q}$.
	Then there exists a unique nef adelic line bundle structure $\bar{L}$ on $L$ such that $f^*\bar{L} \cong \bar{L}^{\otimes q}$.
\end{lemma}
\begin{proof}
	Choose any positive adelic metric $\bar L_0$ on $L$.
	We can choose an open subset $U \subset \Spec \calO_K$ such that $f$ and the isomorphism $f^*L \cong L^{\otimes q}$ extend to $U$.
	Define $\bar L_n \coloneqq q^{-n}(f^n)^*\bar L_0$ on $L$.
	For $v \in U$, the equation follows by the isomorphism.
	Hence we only need to check every $v \notin U$.

	Fix $v \notin U$.
	We have
	\[ \|s(x)\|_n = \|s\circ f^n(x)\|_0^{1/q^n}, \qquad \sup_{x\in X^\an} \left| \log \frac{\|s(f(x))\|_0^{1/q}}{\|s(x)\|_0} \right| \leq C. \]
	Here we view that $s\circ f^n = (f^n)^*s$ as a section of $L^{\otimes q^n}$ by the isomorphism $(f^n)^*L \cong L^{\otimes q^n}$.
	And then
	\[ \dist_v(\bar L_{n+1},\bar L_n) = \frac{1}{q^n} \sup_{y=f^n(x),x\in X^\an} \left| \log \frac{\|s(f(y))\|_0^{1/q}}{\|s(y)\|_0} \right| \leq \frac{C}{q^n}. \]
	Hence $\{\bar L_{n}\}_n$ is a Cauchy sequence and hence converges.

	By the definition of nef adelic metric, the results metric is nef.
	Take limit on the equation $q \bar L_{n+1} \cong f^* \bar L_n$ (note that $f^*$ is continuous), we get $f^*\bar{L} \cong \bar{L}^{\otimes q}$.

	For the uniqueness, choose any other metric $\bar L' = \|\cdots\|'$ on $L$ with $f^*\bar{L'} \cong \bar{L'}^{\otimes q}$.
	Fix a place $v \in M_K$.
	We have
	\[ \dist_v(\bar L',\bar L_n) = \dist_v(q^{-n}(f^n)^*\bar L',\bar L_n) = \frac{1}{q^n} \sup_{y=f^n(x),x\in X^\an} \left| \log \frac{\|s(y)\|'}{\|s(y)\|_0} \right| \leq \frac{C}{q^n} \to 0. \]
	Take limit we get $\bar L = \lim_{n \to \infty} \bar L_n = \bar L'$.
\end{proof}

\begin{definition}[label=\labelNFM{def:canonical-height}, status=draft, tag=0000000]
	Let $X/K$ be a projective variety, $f: X \to X$ a $q$-polarized endomorphism with a $\bar L$ canonical adelic line bundle.
	The \emph{canonical height} of $f$ is defined as the point height $\widehat h_{f,L} \coloneqq h_{\bar L_f}$.
\end{definition}

\begin{proposition}[label=\labelNFM{prop:canonical-height-properties}, status=draft, tag=0000000]
	Let $f:X \to X$ be a $q$-polarized dynamics over $K$ with canonical adelic line bundle $\bar L_f$.
	\begin{enumerate}
		\item The canonical height is nonnegative, and $\forall x \in X(\overline{K})$,
		      \[ \widehat h_{f,L}(x) = 0 \quad\Longleftrightarrow\quad x\text{ is preperiodic}. \]
		\item For any Weil height $h_L$ associated to $L$, one has
		      \[ \widehat h_{f,L}(x) = \lim_{n \to \infty}q^{-n}h_L(f^n(x)), \qquad \widehat h_{f,L} = h_L+O(1). \]
		\item For $X = \bbP^d$, $L = \calO(1)$, and an infinite place $v$, $c_1(\bar L_{f,v}) = T_f$.
	\end{enumerate}
\end{proposition}
\begin{proof}
	The (a) is by the formula $\widehat h(f(x)) = q \widehat h(x)$ and the Northcott property.
	Note that all $f^n(x)$ are defined over $\rkk(x)$.

	Since $q^{-1}h_L(f(x)) - h_L(x)$ is bounded on $X(\overline{K})$, we have
	\[ \left| \frac{h_L(f^{n+1}(x))}{q^{n+1}} - \frac{h_L(f^{n}(x))}{q^{n}} \right| \leq \frac{C}{q^n}. \]
	Hence the sequence is pointwise Cauchy and then pointwise convergence.
	By \Yang{}, $\widehat h$ is also a Weil height associated to $L$, hence they have bounded difference.

	For (c), take $S = c_1(\bar L_{f,v})$.
	It has locally continuous potential $-\log \|s\|_v$.
	Hence by \Yang{}, $q^{-n}(f^n)^*S = S$ converges to $T_f$.
\end{proof}




